% --- Archon physics formalization source begin ---
% archon:physics
% archon:covers PhyXMiniProblems/problem_phyx_mini_0669.lean
% archon:source-report reports/phyx_mini/problem_phyx_mini_0669.source.json

\chapter{Physics problem phyx\_mini\_0669}
\label{ch:PhyXMiniProblems_problem_phyx_mini_0669}

\paragraph{Problem source.}
\#\# Physical scenario

The position versus time for a certain particle moving along the \textbackslash{}( x \textbackslash{}) axis is shown in figure.

\#\# Question

Find the average velocity in the time intervals 0 to \textbackslash{}( 8 \textbackslash{}text\{ s\} \textbackslash{}).

\#\# Answer choices

- A: A: \textbackslash{}[ 6 \textbackslash{}text\{ m/s\} \textbackslash{}]

- B: B: \textbackslash{}[ 10 \textbackslash{}text\{ m/s\} \textbackslash{}]

- C: C: \textbackslash{}[ 0 \textbackslash{}text\{ m/s\} \textbackslash{}]

- D: D: \textbackslash{}[ 2 \textbackslash{}text\{ m/s\} \textbackslash{}]

\#\# Figure caption (auxiliary; use the image as primary evidence)

The image is a line graph illustrating a relation between position \textbackslash{}( x \textbackslash{}) in meters (m) and time \textbackslash{}( t \textbackslash{}) in seconds (s). Here's a detailed description:

- **Axes**:
  - The horizontal axis represents time \textbackslash{}( t \textbackslash{}) in seconds, labeled from 0 to 8.
  - The vertical axis represents position \textbackslash{}( x \textbackslash{}) in meters, labeled from -6 to 10.

- **Data Points and Line**:
  - The graph begins at the origin (0,0).
  - From \textbackslash{}( t = 0 \textbackslash{}) to \textbackslash{}( t = 2 \textbackslash{}), \textbackslash{}( x \textbackslash{}) rises to 10.
  - Between \textbackslash{}( t = 2 \textbackslash{}) and \textbackslash{}( t = 4 \textbackslash{}), \textbackslash{}( x \textbackslash{}) decreases to 5.
  - At \textbackslash{}( t = 5 \textbackslash{}), \textbackslash{}( x \textbackslash{}) remains at 5 until \textbackslash{}( t = 6 \textbackslash{}).
  - From \textbackslash{}( t = 6 \textbackslash{}) to \textbackslash{}( t = 7 \textbackslash{}), \textbackslash{}( x \textbackslash{}) decreases sharply to -5.
  - Between \textbackslash{}( t = 7 \textbackslash{}) and \textbackslash{}( t = 8 \textbackslash{}), \textbackslash{}( x \textbackslash{}) increases back to 0.

- **Grid**:
  - The graph has a grid with horizontal and vertical lines that aid in reading the values.

- **Line Style**:
  - The line is thick and continuous with sharp angles at each data point transition.

The graph displays a fluctuating motion over time, with steady increases, decreases, and periods of constant position.

\#\# Dataset metadata

Kinematics; Spatial Relation Reasoning, Physical Model Grounding Reasoning

\paragraph{Recorded answer/context.}
C: C: \textbackslash{}[ 0 \textbackslash{}text\{ m/s\} \textbackslash{}]

\paragraph{Figure/image path.}
/root/proposal\_for\_physic/science-mango/phyx\_mini\_run/phyx\_data/test\_image/669.png

\paragraph{Formalization target.}
create a compiling Lean file with sorry bodies at `PhyXMiniProblems/problem\_phyx\_mini\_0669.lean`. The Lean declarations must preserve the physical quantities, dimensions or dimensional roles, figure labels, governing-law hypotheses, and final relation expressed by this problem.
Use Mathlib/Physlib names found through LeanExplore where available. If a physics API is missing, introduce faithful local abstractions rather than scalar placeholder aliases.

\begin{theorem}[Physics formalization target]
\label{thm:physics:phyx_mini_0669:target}
\lean{PhyXMiniProblems.ProblemPhyXMini0669.averageVelocityOverDisplayedInterval_is_zero}
\uses{def:physics:phyx-mini-0669:phyxminiproblems-problemphyxmini0669-velocityinmeterspersecond, def:physics:phyx-mini-0669:phyxminiproblems-problemphyxmini0669-matchesprimaryfigure, def:physics:phyx-mini-0669:phyxminiproblems-problemphyxmini0669-particlemotionsetup, def:physics:phyx-mini-0669:phyxminiproblems-problemphyxmini0669-matchesproblemstatement, def:physics:phyx-mini-0669:phyxminiproblems-problemphyxmini0669-satisfiesaveragevelocitylaw, def:physics:phyx-mini-0669:phyxminiproblems-problemphyxmini0669-matchesanswerchoice, def:physics:phyx-mini-0669:phyxminiproblems-problemphyxmini0669-recordedanswerchoice}
The average velocity on the displayed interval is zero metres per second and hence agrees with recorded choice C.
\end{theorem}
\begin{proof}
The horizontal-axis endpoints are $0$ and $8$ seconds, and the graph gives position zero at both endpoints.  Substitute these four readouts into average velocity $[x(8)-x(0)]/(8-0)$ to get zero.  The answer table assigns zero to choice C, and the recorded choice is C, yielding the second conjunct.
\end{proof}

% --- Archon declaration topology begin ---
\section{Declaration topology}
\begin{definition}[Axial Position]
  \label{def:physics:phyx-mini-0669:phyxminiproblems-problemphyxmini0669-axialposition}
  \lean{PhyXMiniProblems.ProblemPhyXMini0669.AxialPosition}
  A signed position coordinate along the particle's one-dimensional axis
\end{definition}
\begin{proof}
  This is the declared model definition of Axial Position.
\end{proof}
\begin{definition}[Axial Velocity]
  \label{def:physics:phyx-mini-0669:phyxminiproblems-problemphyxmini0669-axialvelocity}
  \lean{PhyXMiniProblems.ProblemPhyXMini0669.AxialVelocity}
  A signed one-dimensional velocity, with physical dimension L T⁻¹
\end{definition}
\begin{proof}
  This is the declared model definition of Axial Velocity.
\end{proof}
\begin{definition}[position In Meters]
  \label{def:physics:phyx-mini-0669:phyxminiproblems-problemphyxmini0669-positioninmeters}
  \lean{PhyXMiniProblems.ProblemPhyXMini0669.positionInMeters}
  \uses{def:physics:phyx-mini-0669:phyxminiproblems-problemphyxmini0669-axialposition}
  Metre readout of a signed axial position
\end{definition}
\begin{proof}
  This is the declared model definition of position In Meters.
\end{proof}
\begin{definition}[velocity In Meters Per Second]
  \label{def:physics:phyx-mini-0669:phyxminiproblems-problemphyxmini0669-velocityinmeterspersecond}
  \lean{PhyXMiniProblems.ProblemPhyXMini0669.velocityInMetersPerSecond}
  \uses{def:physics:phyx-mini-0669:phyxminiproblems-problemphyxmini0669-axialvelocity}
  Metres-per-second readout of a signed axial velocity
\end{definition}
\begin{proof}
  This is the declared model definition of velocity In Meters Per Second.
\end{proof}
\begin{definition}[Axis Quantity]
  \label{def:physics:phyx-mini-0669:phyxminiproblems-problemphyxmini0669-axisquantity}
  \lean{PhyXMiniProblems.ProblemPhyXMini0669.AxisQuantity}
  Physical quantity assigned to one of the plot axes
\end{definition}
\begin{proof}
  This is the declared model definition of Axis Quantity.
\end{proof}
\begin{definition}[Axis Unit]
  \label{def:physics:phyx-mini-0669:phyxminiproblems-problemphyxmini0669-axisunit}
  \lean{PhyXMiniProblems.ProblemPhyXMini0669.AxisUnit}
  Unit printed beside a plot axis
\end{definition}
\begin{proof}
  This is the declared model definition of Axis Unit.
\end{proof}
\begin{definition}[Spatial Axis]
  \label{def:physics:phyx-mini-0669:phyxminiproblems-problemphyxmini0669-spatialaxis}
  \lean{PhyXMiniProblems.ProblemPhyXMini0669.SpatialAxis}
  The spatial axis named in the physical scenario
\end{definition}
\begin{proof}
  This is the declared model definition of Spatial Axis.
\end{proof}
\begin{definition}[Trace Geometry]
  \label{def:physics:phyx-mini-0669:phyxminiproblems-problemphyxmini0669-tracegeometry}
  \lean{PhyXMiniProblems.ProblemPhyXMini0669.TraceGeometry}
  The geometry of the thick trace drawn in the primary raster
\end{definition}
\begin{proof}
  This is the declared model definition of Trace Geometry.
\end{proof}
\begin{definition}[Figure Vertex]
  \label{def:physics:phyx-mini-0669:phyxminiproblems-problemphyxmini0669-figurevertex}
  \lean{PhyXMiniProblems.ProblemPhyXMini0669.FigureVertex}
  Names for the six visible vertices of the plotted motion
\end{definition}
\begin{proof}
  This is the declared model definition of Figure Vertex.
\end{proof}
\begin{definition}[Position Time Figure]
  \label{def:physics:phyx-mini-0669:phyxminiproblems-problemphyxmini0669-positiontimefigure}
  \lean{PhyXMiniProblems.ProblemPhyXMini0669.PositionTimeFigure}
  \uses{def:physics:phyx-mini-0669:phyxminiproblems-problemphyxmini0669-axialposition, def:physics:phyx-mini-0669:phyxminiproblems-problemphyxmini0669-axisquantity, def:physics:phyx-mini-0669:phyxminiproblems-problemphyxmini0669-axisunit, def:physics:phyx-mini-0669:phyxminiproblems-problemphyxmini0669-tracegeometry, def:physics:phyx-mini-0669:phyxminiproblems-problemphyxmini0669-figurevertex}
  The calibrated position--time figure. The trace argument is a real scalar because it is the numerical coordinate read from the horizontal axis in seconds; its value is still a dimensionful physical position
\end{definition}
\begin{proof}
  This is the declared model definition of Position Time Figure.
\end{proof}
\begin{definition}[Is Straight Trace Segment]
  \label{def:physics:phyx-mini-0669:phyxminiproblems-problemphyxmini0669-isstraighttracesegment}
  \lean{PhyXMiniProblems.ProblemPhyXMini0669.IsStraightTraceSegment}
  \uses{def:physics:phyx-mini-0669:phyxminiproblems-problemphyxmini0669-positioninmeters, def:physics:phyx-mini-0669:phyxminiproblems-problemphyxmini0669-figurevertex, def:physics:phyx-mini-0669:phyxminiproblems-problemphyxmini0669-positiontimefigure}
  The trace is the straight line segment joining two named consecutive vertices. This is figure evidence, not a dynamical law and not an answer assumption
\end{definition}
\begin{proof}
  This is the declared model definition of Is Straight Trace Segment.
\end{proof}
\begin{definition}[Matches Primary Figure]
  \label{def:physics:phyx-mini-0669:phyxminiproblems-problemphyxmini0669-matchesprimaryfigure}
  \lean{PhyXMiniProblems.ProblemPhyXMini0669.MatchesPrimaryFigure}
  \uses{def:physics:phyx-mini-0669:phyxminiproblems-problemphyxmini0669-positioninmeters, def:physics:phyx-mini-0669:phyxminiproblems-problemphyxmini0669-positiontimefigure, def:physics:phyx-mini-0669:phyxminiproblems-problemphyxmini0669-isstraighttracesegment}
  Exact transcription of image 669. The endpoint positions are calibrated measurements from the graph. They do not state the requested average velocity
\end{definition}
\begin{proof}
  This is the declared model definition of Matches Primary Figure.
\end{proof}
\begin{definition}[Particle Motion Setup]
  \label{def:physics:phyx-mini-0669:phyxminiproblems-problemphyxmini0669-particlemotionsetup}
  \lean{PhyXMiniProblems.ProblemPhyXMini0669.ParticleMotionSetup}
  \uses{def:physics:phyx-mini-0669:phyxminiproblems-problemphyxmini0669-axialvelocity, def:physics:phyx-mini-0669:phyxminiproblems-problemphyxmini0669-spatialaxis, def:physics:phyx-mini-0669:phyxminiproblems-problemphyxmini0669-positiontimefigure}
  The average velocity is an independent physical observable. No numerical value for it is stored in this structure
\end{definition}
\begin{proof}
  This is the declared model definition of Particle Motion Setup.
\end{proof}
\begin{definition}[Matches Problem Statement]
  \label{def:physics:phyx-mini-0669:phyxminiproblems-problemphyxmini0669-matchesproblemstatement}
  \lean{PhyXMiniProblems.ProblemPhyXMini0669.MatchesProblemStatement}
  \uses{def:physics:phyx-mini-0669:phyxminiproblems-problemphyxmini0669-particlemotionsetup}
  The scenario specifies one-dimensional motion along the x axis
\end{definition}
\begin{proof}
  This is the declared model definition of Matches Problem Statement.
\end{proof}
\begin{definition}[Satisfies Average Velocity Law]
  \label{def:physics:phyx-mini-0669:phyxminiproblems-problemphyxmini0669-satisfiesaveragevelocitylaw}
  \lean{PhyXMiniProblems.ProblemPhyXMini0669.SatisfiesAverageVelocityLaw}
  \uses{def:physics:phyx-mini-0669:phyxminiproblems-problemphyxmini0669-positioninmeters, def:physics:phyx-mini-0669:phyxminiproblems-problemphyxmini0669-velocityinmeterspersecond, def:physics:phyx-mini-0669:phyxminiproblems-problemphyxmini0669-particlemotionsetup}
  Average velocity over the displayed interval is displacement divided by its elapsed time. This is the governing kinematic law; it contains no displayed answer value
\end{definition}
\begin{proof}
  This is the declared model definition of Satisfies Average Velocity Law.
\end{proof}
\begin{definition}[Answer Choice]
  \label{def:physics:phyx-mini-0669:phyxminiproblems-problemphyxmini0669-answerchoice}
  \lean{PhyXMiniProblems.ProblemPhyXMini0669.AnswerChoice}
  Labels of the four answer choices printed with the problem
\end{definition}
\begin{proof}
  This is the declared model definition of Answer Choice.
\end{proof}
\begin{definition}[meters Per Second]
  \label{def:physics:phyx-mini-0669:phyxminiproblems-problemphyxmini0669-answerchoice-meterspersecond}
  \lean{PhyXMiniProblems.ProblemPhyXMini0669.AnswerChoice.metersPerSecond}
  \uses{def:physics:phyx-mini-0669:phyxminiproblems-problemphyxmini0669-answerchoice}
  Metres-per-second scalar displayed beside each answer label
\end{definition}
\begin{proof}
  This is the declared model definition of meters Per Second.
\end{proof}
\begin{definition}[Matches Answer Choice]
  \label{def:physics:phyx-mini-0669:phyxminiproblems-problemphyxmini0669-matchesanswerchoice}
  \lean{PhyXMiniProblems.ProblemPhyXMini0669.MatchesAnswerChoice}
  \uses{def:physics:phyx-mini-0669:phyxminiproblems-problemphyxmini0669-axialvelocity, def:physics:phyx-mini-0669:phyxminiproblems-problemphyxmini0669-velocityinmeterspersecond, def:physics:phyx-mini-0669:phyxminiproblems-problemphyxmini0669-answerchoice}
  A physical axial velocity agrees with the scalar printed for a choice
\end{definition}
\begin{proof}
  This is the declared model definition of Matches Answer Choice.
\end{proof}
\begin{definition}[recorded Answer Choice]
  \label{def:physics:phyx-mini-0669:phyxminiproblems-problemphyxmini0669-recordedanswerchoice}
  \lean{PhyXMiniProblems.ProblemPhyXMini0669.recordedAnswerChoice}
  \uses{def:physics:phyx-mini-0669:phyxminiproblems-problemphyxmini0669-answerchoice}
  The answer label recorded in the supplied dataset metadata
\end{definition}
\begin{proof}
  This is the declared model definition of recorded Answer Choice.
\end{proof}
% --- Archon declaration topology end ---
% --- Archon physics formalization source end ---
